\documentclass[11pt]{article}
\usepackage[margin=1in]{geometry}
\usepackage{amsmath,amssymb}
\usepackage{booktabs}
\usepackage[hidelinks]{hyperref}
\usepackage{xcolor}
\setlength{\parindent}{0pt}\setlength{\parskip}{5pt}
%% DRAFT. Target: BlackboxNLP 2026 (EMNLP-colocated), main track, deadline 2026-07-17.
%% Before submission: swap to the BlackboxNLP/ACL style file, anonymize, run the remaining
%% blocking experiments (alignment-null control, multi-layer sweep), and re-check the checklist.
\title{\textbf{Within-Model Redundancy Predicts Co-firing but not Decoder-Geometry
SAE Feature Universality: A Controlled Test of the Natural-Latent Prediction}}
\author{[anonymized for review]}
\date{Draft \today}
\begin{document}
\maketitle

\begin{abstract}
Natural-latent theory predicts that features which are \emph{redundant} within a model (recoverable
from disjoint parts of its activations) should be the ones that recur across independently trained
models. We test this prediction directly on sparse-autoencoder (SAE) features, using a within-model
redundancy score computed from one model's own activations with no access to a second model and no
training. We find a \emph{dissociation} that depends entirely on how ``universality'' is measured.
Under a \emph{decoder-geometry} measure (aligned cross-model decoder cosine), redundancy does
\textbf{not} predict universality on any of three model pairs and is significantly \emph{negative} on
two distant cross-family pairs (partial Spearman $\rho = -0.16$ and $-0.18$); a previously reported
positive on a near pair does not survive a corrected estimator and is attributable to a numerical
artifact. Under a \emph{co-firing} measure (top-context overlap), redundancy shows a positive
association on two of three pairs ($\rho = +0.41$ and $+0.31$), but is null on the third. We conclude
that within-model redundancy tracks \emph{behavioral} feature similarity in some regimes but not
\emph{representational-geometry} similarity, and we caution against using it as a general
cross-model transfer heuristic. All effects are reported as variance-explained with feature-bootstrap
intervals and partials on activation frequency, magnitude, and decoder norm.
\end{abstract}

\section{Introduction}
Feature universality - the observation that independently trained networks learn overlapping features -
is widely reported but under-explained: there is no accepted predictor of \emph{which} features should
transfer. Natural-latent theory offers a candidate. A latent that is \emph{redundant} (recoverable from
each of two disjoint parts of a system) and \emph{mediating} is provably stable across ontologies,
hence a natural target for independently trained models to converge on. This motivates a testable
operational claim: \emph{the SAE features whose activations are most redundant within a model should be
the ones most likely to have a match in a different model.}

We test this claim and report a negative-leaning, metric-dependent result. The contribution is not a new
method but a controlled measurement: (i) under a decoder-geometry notion of universality the redundancy
predictor fails, and a prior positive result is shown to be a harness artifact; (ii) under a co-firing
notion of universality the predictor partially succeeds; (iii) the two conclusions come apart, so the
choice of universality metric - not the redundancy score - drives the headline. This is directly useful
to interpretability researchers testing the natural-abstraction hypothesis and to practitioners who
might otherwise use cheap within-model redundancy as a transfer screen.

\section{Method}
\textbf{Redundancy score.} For each SAE feature $f$ we split the model's own residual activations into
two disjoint halves and measure how well feature $f$'s activation on one half is recoverable, by ridge
regression, from a PCA summary of the other half; the score is the minimum of the two directions,
$\mathrm{redundancy\_R2}(f)=\min(R^2_{\text{left}}, R^2_{\text{right}})$. The regression $R^2$ is floored
at $-1$ (values far below $-1$ indicate a diverged solve, not real un-recoverability); we return to this
in Section~\ref{sec:artifact}. No second model is used.

\textbf{Two universality measures.} Given a second model with its own SAE, we score each feature's
universality two ways. \emph{Decoder-geometry} $U_{\text{dec}}$: align the two decoder spaces and take
the best cross-model decoder cosine, with a shuffled-decoder random baseline. \emph{Co-firing}
$U_{\text{overlap}}$: the Jaccard overlap of the two features' top-activating contexts.

\textbf{Statistics.} We report \emph{partial} Spearman $\rho(\mathrm{redundancy\_R2}, U \mid
\text{frequency}, \text{magnitude}, \text{decoder norm})$ for both measures, a position-shuffle null,
and feature-bootstrap $95\%$ intervals. We report variance-explained ($\rho^2$) as the primary effect
size rather than a $p$-value, given the large feature counts.

\textbf{Model pairs.} A near pair (GPT-2-small vs Pythia-70m-deduped, base models) and two distant
cross-family pairs (Gemma-2-9b vs Llama-3.1-8B; Mistral-7B vs Llama-3.1-8B), each with pretrained SAEs
from different SAE families. All runs use one corrected harness on GPU (Section~\ref{sec:artifact}).

\section{A corrected harness}\label{sec:artifact}
An earlier version of this analysis reported a positive redundancy-universality correlation on the near
pair ($\rho \approx +0.27$ under the decoder measure). That number does not survive two corrections.
First, on large-activation models the ridge recoverability $R^2$ diverged (mean $\approx -1.4\times10^5$)
because the residual PCA basis was ill-conditioned; flooring $R^2$ and unit-variance-scaling the PCA
components restores a sane score (mean redundancy $R^2 \in [0.06, 0.17]$ across pairs). Second, the
decoder-space cross-model alignment must beat a shuffled-decoder baseline; under the corrected,
frequency-partialled estimator the near-pair decoder correlation collapses to $\rho = +0.001$
($p = 0.96$). We therefore treat all pre-correction numbers as artifacts and report only corrected runs.
Every number below is from the corrected harness on GPU (device \texttt{cuda}), with the decoder measure
verified to beat its random baseline.

\section{Results}
\begin{table}[h]\centering
\begin{tabular}{lcc}
\toprule
Model pair & Decoder-geometry $U_{\text{dec}}$ & Co-firing $U_{\text{overlap}}$ \\
 & partial $\rho$ (var.\ expl.) & partial $\rho$ (var.\ expl.) \\
\midrule
GPT-2 / Pythia (near, same-family) & $+0.001$\ \ ($0\%$) & $+0.41$\ \ ($17\%$) \\
Gemma-2-9b / Llama-3.1-8B (far) & $-0.16$\ \ ($2.6\%$) & $+0.02$\ \ ($0\%$) \\
Mistral-7B / Llama-3.1-8B (far) & $-0.18$\ \ ($3.1\%$) & $+0.31$\ \ ($9\%$) \\
\bottomrule
\end{tabular}
\caption{Partial Spearman $\rho$ of within-model redundancy against two universality measures,
controlling for activation frequency, magnitude, and decoder norm. Decoder-geometry: null on the near
pair, significantly negative on both far pairs. Co-firing: strong positive on the near pair and one far
pair, null on the other. All $|\rho|$ small to moderate; see text for bootstrap intervals.}
\end{table}

\textbf{Decoder-geometry universality is not predicted by redundancy} (and is weakly anti-predicted at
distance). The near-pair partial is indistinguishable from zero ($+0.001$), and both distant pairs are
significantly negative ($-0.16$, $-0.18$; each feature-bootstrap $95\%$ interval excludes zero on the
negative side). Redundancy explains at most $\sim$3\% of decoder-universality variance, in the wrong
direction.

\textbf{Co-firing universality is partially predicted by redundancy.} The near pair shows a strong
partial ($+0.41$, $\sim$17\% variance) and the Mistral/Llama far pair a moderate one ($+0.31$,
$\sim$9\%); the Gemma/Llama pair is null ($+0.02$). So the positive is real but pair-dependent, not a
universal law.

\textbf{The dissociation is the finding.} On the same features and the same redundancy score, the sign
and significance of the ``does redundancy predict universality'' question flip with the universality
metric. Redundancy tracks whether two features fire on the \emph{same inputs} (behavioral overlap) in
some regimes, but not whether their decoder directions \emph{align in representation space}.

\section{Discussion}
The natural-latent prediction, taken as a claim about representational geometry, fails here: the
features a model can most redundantly recover from its own activations are not the ones whose decoder
directions match a second model, and at cross-family distance they are slightly less likely to. The
co-firing results show the redundancy signal is not vacuous - it aligns with behavioral similarity - but
behavioral co-firing and geometric alignment are different things, and only the former is (sometimes)
captured. For practitioners, the operational takeaway is negative: within-model redundancy is not a
reliable cross-model transfer screen, especially across architectures.

\section{Limitations}
Three pairs, one layer each, and one redundancy operationalization. The co-firing measure may retain
residual co-activation base-rate structure beyond the frequency partial. The decoder alignment is one
of several possible cross-model matching procedures; an alignment-null control and a multi-layer sweep
are the two most important robustness checks and are in progress. We make no claim about \emph{why} the
Gemma/Llama co-firing result is null while Mistral/Llama is positive; instruct-tuning and SAE-family
differences are candidate explanations we do not resolve.

\section*{Reproducibility}
All results are from a single corrected pipeline; pre-correction numbers are reported only as artifacts
to be superseded. Pair configurations, the $R^2$ floor, the unit-variance PCA conditioning, and the
random-baseline check are specified above.

\end{document}
